\section{Neural Network Architectures for Combinatorial Optimization}

Neural networks have emerged as powerful function approximators for learning solution strategies in combinatorial optimization domains. Unlike traditional approaches that rely on hand-crafted heuristics and domain-specific rules, neural architectures can discover effective decision-making policies by learning patterns from data and experience. This data-driven paradigm enables systems to automatically extract relevant features and decision criteria from problem instances.

The application of neural networks to 3D bin packing involves several distinct architectural paradigms, each offering unique capabilities for processing different aspects of the problem:

\begin{itemize}
\item \textbf{Convolutional Neural Networks (CNNs)}: Process spatial representations such as heightmaps, occupancy grids, and voxelized container states to extract geometric features
\item \textbf{Graph Neural Networks (GNNs)}: Encode relational structure among items, placement locations, and spatial relationships through message passing
\item \textbf{Recurrent Neural Networks (RNNs)}: Model sequential dependencies in packing trajectories and temporal evolution of container states
\item \textbf{Reinforcement Learning Architectures}: Learn placement policies through trial-and-error interaction with packing environments (e.g., Deep Q-Networks)
\item \textbf{Attention Mechanisms}: Enable relational reasoning over variable-size sets of items and candidate positions
\end{itemize}

This section reviews these architectural families, examining how each provides distinct inductive biases for capturing problem structure. We begin with foundational architectures for spatial and relational processing, then explore reinforcement learning frameworks that leverage these architectures for sequential decision-making in packing problems.


\subsection{Convolutional Neural Networks (CNNs)}

Convolutional Neural Networks excel at processing grid-structured spatial data, making them particularly well-suited for bin packing representations that discretize the container space. CNNs automatically learn hierarchical spatial features through local receptive fields and weight sharing, enabling efficient processing of heightmaps, occupancy grids, and voxelized container states.

\subsubsection{Architectural Components}

A typical CNN architecture for packing consists of several key components:

\paragraph{Convolutional Layers:}

The fundamental operation is the 2D or 3D convolution:

\begin{equation}
(\mathbf{f} * \mathbf{w})[i, j, k] = \sum_{m,n,p} \mathbf{f}[i+m, j+n, k+p] \cdot \mathbf{w}[m, n, p]
\end{equation}

where $\mathbf{f}$ is the input feature map (e.g., heightmap or occupancy grid), $\mathbf{w}$ is a learnable filter kernel, and the output represents a detected spatial pattern. Multiple filters are applied in parallel to detect diverse features (edges, corners, filled regions, empty spaces).

\paragraph{Receptive Fields and Hierarchy:}

CNNs build hierarchical representations through stacked layers:
\begin{itemize}
\item \textbf{Early layers}: Detect local patterns (edges, corners, small obstacles)
\item \textbf{Middle layers}: Combine local patterns into larger structures (walls, clusters of items, empty regions)
\item \textbf{Deep layers}: Capture global context (overall container utilization, large-scale spatial organization)
\end{itemize}

\paragraph{Pooling Operations:}

Pooling layers reduce spatial dimensions while retaining important features:

\begin{itemize}
\item \textbf{Max pooling}: $\text{MaxPool}(\mathbf{f})[i, j] = \max_{m,n \in W} \mathbf{f}[i+m, j+n]$
  \\ Preserves strong activations (e.g., presence of items)
\item \textbf{Average pooling}: Captures aggregate statistics over regions
\end{itemize}

\subsubsection{Spatial Representations for Packing}

CNNs process various spatial encodings of the packing state:

\paragraph{Heightmaps (2D):}

A 2D grid where each cell $(x, y)$ stores the height of items at that position:

\begin{equation}
H[x, y] = \max_{z} \{\text{item occupies }(x, y, z)\}
\end{equation}

\begin{itemize}
\item \textbf{Advantages}: Compact representation, fast to update, captures vertical structure
\item \textbf{Limitations}: Loses information about internal voids and overhangs
\end{itemize}

\paragraph{Occupancy Grids (2D slices or 3D voxels):}

Binary or multi-valued grids indicating whether each cell is occupied:

\begin{equation}
O[x, y, z] = \begin{cases}
1 & \text{if cell occupied} \\
0 & \text{if cell empty}
\end{cases}
\end{equation}

\begin{itemize}
\item \textbf{3D voxel grids}: Complete geometric representation but memory-intensive ($O(W \times D \times H)$)
\item \textbf{2D slices}: Horizontal or vertical cross-sections provide partial 3D information with lower memory cost
\end{itemize}

\paragraph{Multi-Channel Representations:}

Combine multiple spatial features as channels (analogous to RGB channels in images):

\begin{itemize}
\item Channel 1: Binary occupancy
\item Channel 2: Item IDs or types
\item Channel 3: Height values
\item Channel 4: Stability indicators
\item Channel 5: Distance to container boundaries
\end{itemize}

\subsubsection{Application to Placement Decisions}

CNNs can be integrated into packing systems in several ways:

\paragraph{1. Candidate Position Scoring:}

Given a candidate placement location $(x, y, z)$ for the current item:

\begin{enumerate}
\item \textbf{Local patch extraction}: Extract a spatial patch centered on the candidate position (e.g., $32 \times 32 \times 32$ voxel region)
\item \textbf{CNN processing}: Pass the patch through convolutional layers to extract spatial features
\item \textbf{Placement quality prediction}: Final fully-connected layers output a score indicating placement quality
\end{enumerate}

\paragraph{2. Global State Encoding:}

Process the entire container state to extract global features:

\begin{enumerate}
\item \textbf{Full container representation}: Heightmap or occupancy grid of the entire container
\item \textbf{Convolutional processing}: Multiple conv layers reduce spatial dimensions while increasing feature depth
\item \textbf{Global feature vector}: Flatten and concatenate with item features
\item \textbf{Decision network}: Use global features to score all candidate positions
\end{enumerate}

\paragraph{3. Hybrid Spatial-Vector Approach:}

\begin{itemize}
\item \textbf{CNN branch}: Processes spatial representation $\to$ spatial features $\mathbf{f}_{\text{spatial}} \in \mathbb{R}^{d_s}$
\item \textbf{MLP branch}: Processes item and candidate features $\to$ semantic features $\mathbf{f}_{\text{semantic}} \in \mathbb{R}^{d_m}$
\item \textbf{Fusion}: Concatenate $[\mathbf{f}_{\text{spatial}}; \mathbf{f}_{\text{semantic}}]$ and feed to decision head
\end{itemize}

\subsubsection{Architectural Variations}

\paragraph{ResNet-style Skip Connections:}

For deeper networks, residual connections prevent gradient vanishing:

\begin{equation}
\mathbf{h}^{(\ell+1)} = \sigma(\text{Conv}(\mathbf{h}^{(\ell)})) + \mathbf{h}^{(\ell)}
\end{equation}

This enables training of very deep CNNs (50+ layers) that can capture complex spatial patterns.

\paragraph{Dilated Convolutions:}

Expand receptive fields without increasing parameters:

\begin{equation}
(\mathbf{f} *_d \mathbf{w})[i] = \sum_{m} \mathbf{f}[i + d \cdot m] \cdot \mathbf{w}[m]
\end{equation}

where $d$ is the dilation rate. Useful for capturing long-range spatial dependencies (e.g., items at opposite ends of the container).

\paragraph{3D Convolutions:}

For true volumetric processing:

\begin{itemize}
\item \textbf{Filters}: $\mathbf{w} \in \mathbb{R}^{k \times k \times k \times c_{\text{in}} \times c_{\text{out}}}$
\item \textbf{Operations}: Convolve over all three spatial dimensions
\item \textbf{Cost}: Significantly higher computational and memory requirements
\end{itemize}

\subsubsection{Advantages for Packing Problems}

\begin{itemize}
\item \textbf{Translation equivariance}: Same features are detected regardless of position in the container
\item \textbf{Parameter efficiency}: Weight sharing dramatically reduces parameters compared to fully-connected networks
\item \textbf{Hierarchical feature learning}: Automatically learns relevant spatial patterns from data
\item \textbf{Natural fit for spatial domains}: Directly processes geometric representations without manual feature engineering
\end{itemize}

\subsubsection{Limitations}

\begin{itemize}
\item \textbf{Discretization requirement}: Continuous 3D space must be voxelized, leading to precision loss
\item \textbf{Memory scaling}: 3D voxel grids scale cubically with container dimensions ($O(W \times D \times H)$)
\item \textbf{Fixed-size inputs}: Standard CNNs require fixed spatial dimensions (addressed by adaptive pooling or fully-convolutional architectures)
\item \textbf{Limited relational reasoning}: Primarily captures local spatial patterns; struggles with long-range dependencies between distant items
\end{itemize}


\subsection{Graph Neural Networks (GNNs)}

Graph Neural Networks provide a natural framework for encoding the \textbf{relational structure} inherent in packing problems: items, placement locations, and spatial relationships can be explicitly represented as nodes and edges in a graph \cite{SZL+22, ZWL+24}. Unlike CNNs that impose a grid structure, GNNs work directly with arbitrary graph topologies, making them ideal for capturing pairwise relationships and constraints.

\subsubsection{Graph Representation of Packing State}

A packing configuration can be modeled as a graph $G = (V, E)$ where:

\paragraph{Nodes $V$:}
\begin{itemize}
\item \textbf{Placed item nodes}: Each placed item $i$ with features encoding its geometry and position
\item \textbf{Empty space nodes}: Available placement locations with features describing size and position
\item \textbf{Container node}: Global state (utilization, capacity, packing progress)
\item \textbf{Current item node}: The item currently being placed
\end{itemize}

\paragraph{Edges $E$:}
\begin{itemize}
\item \textbf{Support edges}: Connections between items in physical contact
\item \textbf{Proximity edges}: Connections between nearby items or spaces
\item \textbf{Feasibility edges}: Connections indicating valid placement options
\item \textbf{Global edges}: Connections to container node for global context
\end{itemize}

\subsubsection{Message Passing Framework}

GNNs learn node representations by iteratively \textbf{aggregating information} from neighbors. At each layer $\ell$:

\begin{enumerate}
\item \textbf{Message computation}: For each edge $(i, j)$:
   \begin{equation}
   \mathbf{m}_{ij}^{(\ell)} = \text{Message}(\mathbf{h}_i^{(\ell-1)}, \mathbf{h}_j^{(\ell-1)}, \mathbf{e}_{ij})
   \end{equation}
   where $\mathbf{h}_i$ is the hidden state of node $i$, and $\mathbf{e}_{ij}$ are edge features (e.g., distance, contact area).

\item \textbf{Aggregation}: Collect messages from all neighbors:
   \begin{equation}
   \mathbf{a}_i^{(\ell)} = \text{Aggregate}\left(\{\mathbf{m}_{ji}^{(\ell)} : j \in \mathcal{N}(i)\}\right)
   \end{equation}
   Common aggregators: sum, mean, max, or attention-weighted sum.

\item \textbf{Update}: Combine aggregated messages with current state:
   \begin{equation}
   \mathbf{h}_i^{(\ell)} = \text{Update}(\mathbf{h}_i^{(\ell-1)}, \mathbf{a}_i^{(\ell)})
   \end{equation}
   Typically a neural network (MLP or GRU).
\end{enumerate}

After $L$ layers, each node has a representation $\mathbf{h}_i^{(L)}$ that aggregates information from its $L$-hop neighborhood.

\subsubsection{GNN Variants for Packing}

\paragraph{Graph Convolutional Networks (GCN) \cite{KW17}:}
\begin{equation}
\mathbf{h}_i^{(\ell)} = \sigma\left(W^{(\ell)} \sum_{j \in \mathcal{N}(i) \cup \{i\}} \frac{\mathbf{h}_j^{(\ell-1)}}{\sqrt{|\mathcal{N}(i)||\mathcal{N}(j)|}}\right)
\end{equation}

Simple and efficient, but fixed aggregation (mean).

\paragraph{Graph Attention Networks (GAT) \cite{VCC+18}:}
\begin{equation}
\mathbf{h}_i^{(\ell)} = \sigma\left(\sum_{j \in \mathcal{N}(i)} \alpha_{ij} W^{(\ell)} \mathbf{h}_j^{(\ell-1)}\right)
\end{equation}

where $\alpha_{ij}$ are learned attention weights. Allows the network to focus on relevant neighbors (e.g., items that directly support the current item).

\paragraph{Relational GCN (R-GCN) \cite{SKB+18}:}

Uses different weight matrices for different edge types:
\begin{equation}
\mathbf{h}_i^{(\ell)} = \sigma\left(\sum_{r \in \mathcal{R}} \sum_{j \in \mathcal{N}_r(i)} \frac{1}{|\mathcal{N}_r(i)|} W_r^{(\ell)} \mathbf{h}_j^{(\ell-1)}\right)
\end{equation}

where $\mathcal{R}$ is the set of relation types (support, proximity, containment). Useful when different relationships have different semantics.

\subsubsection{Application to Placement Selection}

Given a graph representation of the current packing state:

\begin{enumerate}
\item \textbf{Graph construction}: Build $G$ from placed items, EMS candidates, and the item to place
\item \textbf{Message passing}: Run $L$ GNN layers to compute node embeddings $\{\mathbf{h}_i^{(L)}\}$
\item \textbf{Readout}:
  \begin{itemize}
  \item \textbf{Node-level}: Score each EMS node: $Q(s, a_j) = \text{MLP}(\mathbf{h}_j^{(L)})$
  \item \textbf{Graph-level}: Pool all EMS embeddings and decode: $Q(s, a) = \text{MLP}(\text{Pool}(\{\mathbf{h}_j^{(L)}\}))$
  \end{itemize}
\end{enumerate}

\subsubsection{Advantages for Packing Problems}

\begin{itemize}
\item \textbf{Explicit relational structure}: Encodes support, adjacency, and blocking relationships directly
\item \textbf{Inductive bias}: Message passing naturally captures spatial dependencies (e.g., ``supported items depend on supporters'')
\item \textbf{Interpretability}: Attention weights or message values reveal which relationships drive decisions
\item \textbf{Variable-size graphs}: Handles varying numbers of items and placement candidates without architectural changes
\end{itemize}

\subsubsection{Challenges}

\begin{itemize}
\item \textbf{Graph construction overhead}: Building and updating the graph at each packing step is computationally expensive
\item \textbf{Dynamic graphs}: Adding/removing nodes (items/EMS) requires re-indexing and re-initializing embeddings
\item \textbf{Scalability}: Large packing problems (100s of items) lead to dense graphs with $O(n^2)$ edges
\item \textbf{Limited adoption}: GNNs for bin packing remain less common than CNNs or Transformers, with fewer established architectures
\end{itemize}


\subsection{Recurrent Neural Networks (RNNs)}

Recurrent Neural Networks excel at processing \textbf{sequential data} by maintaining an internal hidden state that evolves as the network processes elements one at a time. In packing problems, RNNs can model the temporal dependencies in placement sequences: each placement decision affects the container state and constrains future options, creating a sequential decision trajectory.

\subsubsection{Core RNN Architecture}

The fundamental RNN update equations are:

\begin{align}
\mathbf{h}_t &= \sigma(\mathbf{W}_{hh} \mathbf{h}_{t-1} + \mathbf{W}_{xh} \mathbf{x}_t + \mathbf{b}_h) \\
\mathbf{y}_t &= \mathbf{W}_{hy} \mathbf{h}_t + \mathbf{b}_y
\end{align}

where:
\begin{itemize}
\item $\mathbf{x}_t$: Input at time step $t$ (e.g., current item and container features)
\item $\mathbf{h}_t$: Hidden state capturing history of previous placements
\item $\mathbf{y}_t$: Output (e.g., placement decision or value estimate)
\item $\mathbf{W}_{hh}, \mathbf{W}_{xh}, \mathbf{W}_{hy}$: Learnable weight matrices
\end{itemize}

The hidden state $\mathbf{h}_t$ serves as a \textbf{memory} of the packing trajectory, encoding which items have been placed and how the container state has evolved.

\subsubsection{Long Short-Term Memory (LSTM)}

Standard RNNs suffer from \textbf{vanishing gradients}, making it difficult to learn long-range dependencies. LSTMs address this through a gating mechanism that selectively preserves and updates information:

\begin{align}
\mathbf{f}_t &= \sigma(\mathbf{W}_f [\mathbf{h}_{t-1}, \mathbf{x}_t] + \mathbf{b}_f) \quad \text{(forget gate)} \\
\mathbf{i}_t &= \sigma(\mathbf{W}_i [\mathbf{h}_{t-1}, \mathbf{x}_t] + \mathbf{b}_i) \quad \text{(input gate)} \\
\tilde{\mathbf{c}}_t &= \tanh(\mathbf{W}_c [\mathbf{h}_{t-1}, \mathbf{x}_t] + \mathbf{b}_c) \quad \text{(candidate cell state)} \\
\mathbf{c}_t &= \mathbf{f}_t \odot \mathbf{c}_{t-1} + \mathbf{i}_t \odot \tilde{\mathbf{c}}_t \quad \text{(cell state update)} \\
\mathbf{o}_t &= \sigma(\mathbf{W}_o [\mathbf{h}_{t-1}, \mathbf{x}_t] + \mathbf{b}_o) \quad \text{(output gate)} \\
\mathbf{h}_t &= \mathbf{o}_t \odot \tanh(\mathbf{c}_t) \quad \text{(hidden state)}
\end{align}

where $\odot$ denotes element-wise multiplication. The gates control:
\begin{itemize}
\item \textbf{Forget gate $\mathbf{f}_t$}: What to discard from previous cell state (e.g., forget items from early placements)
\item \textbf{Input gate $\mathbf{i}_t$}: What new information to store (e.g., update container utilization)
\item \textbf{Output gate $\mathbf{o}_t$}: What to expose from the cell state
\end{itemize}

The cell state $\mathbf{c}_t$ acts as a \textbf{long-term memory} that can propagate information across many time steps.

\subsubsection{Gated Recurrent Units (GRU)}

GRUs simplify the LSTM architecture while retaining similar performance:

\begin{align}
\mathbf{z}_t &= \sigma(\mathbf{W}_z [\mathbf{h}_{t-1}, \mathbf{x}_t]) \quad \text{(update gate)} \\
\mathbf{r}_t &= \sigma(\mathbf{W}_r [\mathbf{h}_{t-1}, \mathbf{x}_t]) \quad \text{(reset gate)} \\
\tilde{\mathbf{h}}_t &= \tanh(\mathbf{W}_h [\mathbf{r}_t \odot \mathbf{h}_{t-1}, \mathbf{x}_t]) \quad \text{(candidate hidden state)} \\
\mathbf{h}_t &= (1 - \mathbf{z}_t) \odot \mathbf{h}_{t-1} + \mathbf{z}_t \odot \tilde{\mathbf{h}}_t \quad \text{(hidden state)}
\end{align}

GRUs have fewer parameters than LSTMs (no separate cell state), making them faster to train while still handling long-range dependencies.

\subsubsection{Application to Packing Sequences}

RNNs model the packing process as a sequence of decisions:

\paragraph{1. Sequential Item Encoding:}

Process the item sequence $\{\mathbf{v}_1, \mathbf{v}_2, \ldots, \mathbf{v}_n\}$ through an RNN:

\begin{equation}
\mathbf{h}_i = \text{RNN}(\mathbf{v}_i, \mathbf{h}_{i-1})
\end{equation}

The final hidden state $\mathbf{h}_n$ encodes the entire item set, capturing ordering and composition.

\paragraph{2. Packing Trajectory Modeling:}

At each placement step $t$:
\begin{enumerate}
\item \textbf{Input}: Current container state features and item to place: $\mathbf{x}_t = [\mathbf{s}_{\text{container}}, \mathbf{v}_{\text{current}}]$
\item \textbf{Hidden state update}: $\mathbf{h}_t = \text{LSTM}(\mathbf{x}_t, \mathbf{h}_{t-1})$
\item \textbf{Decision}: Use $\mathbf{h}_t$ to score candidate placements or directly output a placement decision
\end{enumerate}

The hidden state $\mathbf{h}_t$ captures:
\begin{itemize}
\item \textbf{Placement history}: Which items have been placed and where
\item \textbf{Container evolution}: How utilization and available space have changed
\item \textbf{Sequential patterns}: Regularities in successful packing trajectories
\end{itemize}

\paragraph{3. Encoder-Decoder Architecture:}

\begin{itemize}
\item \textbf{Encoder RNN}: Processes all items to produce a context vector $\mathbf{c} = \mathbf{h}_n$
\item \textbf{Decoder RNN}: Generates placement sequence conditioned on $\mathbf{c}$:
  \begin{equation}
  \mathbf{h}_t^{\text{dec}} = \text{RNN}(\mathbf{c}, \mathbf{h}_{t-1}^{\text{dec}})
  \end{equation}
\item \textbf{Attention mechanism}: Decoder can attend to encoder hidden states to focus on relevant items
\end{itemize}

\paragraph{4. Bidirectional RNNs:}

For offline packing (all items known in advance), process the item sequence in both directions:

\begin{align}
\overrightarrow{\mathbf{h}}_i &= \text{RNN}(\mathbf{v}_i, \overrightarrow{\mathbf{h}}_{i-1}) \quad \text{(forward)} \\
\overleftarrow{\mathbf{h}}_i &= \text{RNN}(\mathbf{v}_i, \overleftarrow{\mathbf{h}}_{i+1}) \quad \text{(backward)} \\
\mathbf{h}_i &= [\overrightarrow{\mathbf{h}}_i; \overleftarrow{\mathbf{h}}_i] \quad \text{(concatenate)}
\end{align}

This captures both past and future context for each item, improving placement decisions.

\subsubsection{Advantages for Packing Problems}

\begin{itemize}
\item \textbf{Temporal modeling}: Naturally captures sequential dependencies in packing trajectories
\item \textbf{Variable-length sequences}: Handles varying numbers of items without architectural changes
\item \textbf{Memory of history}: Hidden state maintains context of previous placements
\item \textbf{Order-aware}: Explicitly models the importance of item ordering (unlike order-invariant architectures)
\end{itemize}

\subsubsection{Limitations}

\begin{itemize}
\item \textbf{Sequential processing}: Cannot parallelize over time steps, slower than Transformers for long sequences
\item \textbf{Limited long-range dependencies}: Despite LSTM/GRU improvements, still struggles with very long sequences (100+ items)
\item \textbf{Order sensitivity}: Performance depends on item ordering, may need to learn ordering jointly
\item \textbf{Gradient instability}: Training deep RNNs can be challenging, especially for long sequences
\end{itemize}

\subsubsection{Comparison with Transformers}

While both RNNs and Transformers handle sequences, they differ fundamentally:

\begin{table}[h]
\centering
\begin{tabular}{|l|p{5.5cm}|p{5.5cm}|}
\hline
\textbf{Property} & \textbf{RNN/LSTM} & \textbf{Transformer} \\
\hline
Processing & Sequential (one step at a time) & Parallel (all positions simultaneously) \\
\hline
Position encoding & Implicit (via recurrence) & Explicit (positional embeddings) \\
\hline
Long-range dependencies & Limited (vanishing gradients) & Excellent (direct attention) \\
\hline
Computational cost & $O(n \cdot d^2)$ & $O(n^2 \cdot d)$ \\
\hline
Training speed & Slower (sequential) & Faster (parallelizable) \\
\hline
Inference speed & Fast for short sequences & Fast for any length \\
\hline
\end{tabular}
\end{table}


\subsection{Deep Q-Networks (DQN)}

Deep Q-Networks revolutionized reinforcement learning by combining Q-learning with deep neural networks, enabling agents to learn policies directly from high-dimensional state representations \cite{MKS+15}. DQN provides a natural framework for sequential decision problems with discrete action spaces: the agent learns to evaluate actions based on their long-term value, building up policies that maximize cumulative reward over entire solution trajectories.

\subsubsection{The Q-Learning Foundation}

Q-learning seeks to learn an action-value function $Q(s, a)$ that estimates the expected cumulative reward of taking action $a$ in state $s$ and following the optimal policy thereafter. The optimal Q-function satisfies the \textbf{Bellman equation}:

\begin{equation}
Q^*(s, a) = \mathbb{E}_{s'}\left[r + \gamma \max_{a'} Q^*(s', a')\right]
\end{equation}

where $r$ is the immediate reward, $\gamma \in [0, 1]$ is the discount factor, and $s'$ is the next state. Traditional Q-learning uses a table to store $Q(s, a)$ for every state-action pair, which is infeasible for large or continuous state spaces.

\subsubsection{Neural Function Approximation}

DQN replaces the Q-table with a neural network $Q(s, a; \theta)$ parameterized by weights $\theta$. The network is trained to minimize the \textbf{temporal difference (TD) error}:

\begin{equation}
\mathcal{L}(\theta) = \mathbb{E}_{(s, a, r, s') \sim \mathcal{D}}\left[\left(r + \gamma \max_{a'} Q(s', a'; \theta^-) - Q(s, a; \theta)\right)^2\right]
\end{equation}

where $\mathcal{D}$ is a \textbf{replay buffer} storing past experiences $(s, a, r, s')$, and $\theta^-$ are the parameters of a \textbf{target network} that is updated periodically to stabilize training.

\paragraph{Key Innovations:}

\begin{enumerate}
\item \textbf{Experience Replay}: Past transitions $(s, a, r, s')$ are stored in a buffer and randomly sampled during training. This breaks temporal correlations in sequential data and improves sample efficiency.

\item \textbf{Target Network}: A separate network with frozen weights $\theta^-$ is used to compute the TD target $r + \gamma \max_{a'} Q(s', a'; \theta^-)$. The target network is updated every $C$ steps by copying $\theta \to \theta^-$, preventing divergence caused by the ``moving target'' problem.

\item \textbf{$\varepsilon$-greedy Exploration}: The agent selects the greedy action $a = \arg\max_a Q(s, a; \theta)$ with probability $1 - \varepsilon$, and a random action with probability $\varepsilon$. Epsilon is annealed from $\varepsilon_{\text{start}}$ to $\varepsilon_{\text{end}}$ over training.
\end{enumerate}

\subsubsection{Application to Packing Problems}

In packing domains, the DQN framework is instantiated as follows:

\begin{itemize}
\item \textbf{State $s$}: The current packing configuration, typically encoded as:
  \begin{itemize}
  \item Feature vectors capturing item and container properties (dimensions, weight, shape)
  \item Spatial representations (heightmaps, occupancy grids, voxel representations)
  \item Global statistics (utilization, number of items placed/remaining)
  \item Current item awaiting placement
  \end{itemize}

\item \textbf{Action $a$}: A discrete choice among placement candidates. For $K$ candidate positions (e.g., corner points, empty maximal spaces), $a \in \{1, \ldots, K\}$. Actions may also encode orientation selection or bin selection in multi-container scenarios.

\item \textbf{Reward $r$}: Immediate feedback signal that guides learning. Common reward structures include:
  \begin{itemize}
  \item \textbf{Utilization increment}: Reward proportional to volume successfully packed
  \item \textbf{Sparse completion reward}: Reward only at episode end based on final packing quality
  \item \textbf{Penalty-based}: Negative reward for opening new containers or violating constraints
  \item \textbf{Shaped rewards}: Additional terms encouraging desirable properties (stability, compactness, accessibility)
  \end{itemize}

\item \textbf{Episode Termination}: When all items are successfully placed or no feasible placements remain.
\end{itemize}

\paragraph{Network Architecture:}

The Q-network typically consists of:

\begin{enumerate}
\item \textbf{Input encoding}: State features $\mathbf{s} \in \mathbb{R}^d$ (item + container + global features)
\item \textbf{Shared trunk}: Multi-layer perceptron (MLP) or convolutional layers to process spatial representations
\item \textbf{Action-value heads}: Outputs $Q(s, a)$ for each candidate action $a$
\end{enumerate}

For $K$ candidate placements, one can use:
\begin{itemize}
\item \textbf{Separate forward passes}: Encode each $(state, candidate_k)$ pair and evaluate $Q(s, a_k)$ separately
\item \textbf{Single forward pass}: Encode state once, then score all candidates via attention or separate heads
\end{itemize}

\paragraph{Advantages:}

\begin{itemize}
\item \textbf{End-to-end learning}: No need to manually design scoring functions; the network learns what features correlate with good outcomes
\item \textbf{Generalization}: Networks trained on diverse problem instances can generalize to unseen items and container sizes
\item \textbf{Implicit constraint handling}: Through negative rewards, the agent learns to avoid infeasible placements
\end{itemize}

\paragraph{Limitations:}

\begin{itemize}
\item \textbf{Overestimation bias}: The $\max$ operator in the Bellman update introduces positive bias, leading to overoptimistic Q-values and unstable learning (addressed by Double DQN)
\item \textbf{Discrete action spaces}: DQN requires discretization of continuous placement positions, limiting precision
\item \textbf{Sample inefficiency}: Requires many episodes to learn effective policies, especially for large combinatorial action spaces
\item \textbf{Hyperparameter sensitivity}: Learning rate, replay buffer size, target update frequency, and epsilon schedule require careful tuning
\end{itemize}


\subsection{Double DQN}

Standard DQN suffers from \textbf{overestimation bias}: the $\max$ operator in the Bellman target tends to select overestimated action values, leading to inflated Q-values and unstable learning \cite{HGS16}. Double DQN addresses this by decoupling action selection from action evaluation.

\subsubsection{The Overestimation Problem}

In DQN, the target is computed as:

\begin{equation}
y_{\text{DQN}} = r + \gamma \max_{a'} Q(s', a'; \theta^-)
\end{equation}

The same network (or target network) is used both to \textbf{select} the best action ($\arg\max$) and \textbf{evaluate} its value ($Q$). If the Q-function has estimation errors, the $\max$ will preferentially select actions with \textbf{positive errors}, causing systematic overestimation.

\subsubsection{Double Q-Learning Solution}

Double DQN uses the \textbf{online network} to select the action and the \textbf{target network} to evaluate it:

\begin{equation}
y_{\text{Double}} = r + \gamma Q(s', \arg\max_{a'} Q(s', a'; \theta); \theta^-)
\end{equation}

\paragraph{Breakdown:}

\begin{enumerate}
\item \textbf{Action selection}: Use current weights $\theta$ to find the best action:
   \begin{equation}
   a^* = \arg\max_{a'} Q(s', a'; \theta)
   \end{equation}

\item \textbf{Action evaluation}: Use target weights $\theta^-$ to evaluate that action:
   \begin{equation}
   Q(s', a^*; \theta^-)
   \end{equation}
\end{enumerate}

This separation reduces overestimation because positive errors in the online network (selecting $a^*$) are unlikely to correlate with positive errors in the target network (evaluating $a^*$).

\subsubsection{Impact on Packing Problems}

For packing applications, Double DQN provides:

\begin{itemize}
\item \textbf{More stable Q-values}: Especially important when reward signals are sparse (e.g., only at episode end)
\item \textbf{Better convergence}: Reduced oscillation in learned policies during training
\item \textbf{Improved sample efficiency}: Faster learning with fewer episodes
\end{itemize}

Empirical results \cite{HGS16} show that Double DQN often achieves higher final performance and more consistent learning curves compared to standard DQN, particularly in domains with large action spaces---such as selecting among many candidate placement positions when numerous items remain to be packed.


\subsection{Transformer Architectures}

Transformers, originally developed for natural language processing \cite{VSP+17}, have become a dominant architecture for modeling relational structures in combinatorial optimization \cite{KW19, XGP+24}. Unlike convolutional or recurrent networks, Transformers process \textbf{sets} of elements (items, placements, spatial regions) via \textbf{self-attention}, learning which relationships matter for decision-making.

\subsubsection{The Attention Mechanism}

The core operation is \textbf{scaled dot-product attention}:

\begin{equation}
\text{Attention}(Q, K, V) = \text{softmax}\left(\frac{QK^\top}{\sqrt{d_k}}\right) V
\end{equation}

where:
\begin{itemize}
\item $Q \in \mathbb{R}^{n \times d_k}$: \textbf{Queries} (what we're looking for)
\item $K \in \mathbb{R}^{m \times d_k}$: \textbf{Keys} (what's available)
\item $V \in \mathbb{R}^{m \times d_v}$: \textbf{Values} (content to aggregate)
\item $d_k$: Dimension of keys/queries (scaling prevents gradient vanishing)
\end{itemize}

\textbf{Interpretation}: For each query $q_i$, compute attention weights over all keys $k_j$ via softmax of dot products, then aggregate the corresponding values $v_j$. High $q_i \cdot k_j$ means element $i$ ``attends to'' element $j$.

\subsubsection{Multi-Head Attention}

To capture diverse relationships, Transformers use \textbf{multiple attention heads}:

\begin{equation}
\text{MultiHead}(Q, K, V) = \text{Concat}(\text{head}_1, \ldots, \text{head}_h) W^O
\end{equation}

where each head learns different attention patterns:

\begin{equation}
\text{head}_i = \text{Attention}(QW_i^Q, KW_i^K, VW_i^V)
\end{equation}

Different heads might focus on:
\begin{itemize}
\item Spatial proximity (items near each other)
\item Size compatibility (items with similar dimensions)
\item Constraint relationships (fragile items avoiding heavy items above)
\end{itemize}

\subsubsection{Transformer Encoder}

A full Transformer encoder layer consists of:

\begin{enumerate}
\item \textbf{Multi-head self-attention}: Elements attend to each other
\item \textbf{Feed-forward network}: Per-element MLP
\item \textbf{Residual connections} and \textbf{layer normalization} for stable training
\end{enumerate}

\begin{align}
\mathbf{Z} &= \text{LayerNorm}(\mathbf{X} + \text{MultiHead}(\mathbf{X}, \mathbf{X}, \mathbf{X})) \\
\mathbf{H} &= \text{LayerNorm}(\mathbf{Z} + \text{FFN}(\mathbf{Z}))
\end{align}

Multiple encoder layers are stacked to form deep representations.

\subsubsection{Application to Packing Problems}

Transformers are well-suited to packing because these problems fundamentally involve \textbf{relational reasoning} over sets:

\begin{itemize}
\item \textbf{Items}: Which items should be placed together? Which have compatible shapes and sizes?
\item \textbf{Placements}: Which candidate positions create stable, efficient configurations?
\item \textbf{Constraints}: How do current placements affect future options?
\end{itemize}

\paragraph{Example: Packing Transformers}

Recent work has applied Transformer architectures to packing domains, demonstrating their effectiveness:

\begin{enumerate}
\item \textbf{Item Encoding}: Each item $i$ is embedded as $\mathbf{e}_i \in \mathbb{R}^d$ capturing relevant properties (dimensions, weight, geometric features)

\item \textbf{Candidate Position Encoding}: Each candidate placement location $j$ is embedded as $\mathbf{c}_j \in \mathbb{R}^d$ (spatial coordinates, available space, local geometry)

\item \textbf{Cross-Attention}: Items attend to placement candidates:
   \begin{equation}
   \mathbf{h}_j = \text{Attention}(Q=\mathbf{c}_j, K=\{\mathbf{e}_i\}, V=\{\mathbf{e}_i\})
   \end{equation}

   This computes, for each candidate position $j$, which items are relevant based on compatibility and packing sequence.

\item \textbf{Spatial Feature Integration}: Spatial representations (e.g., heightmaps, occupancy grids) may be processed by convolutional layers to extract spatial features $\mathbf{f}_{\text{spatial}}$, then concatenated with item and position embeddings.

\item \textbf{Decision Head}: The attended representations are passed to output layers that produce placement probabilities or value estimates.
\end{enumerate}

\paragraph{Benefits for Packing:}

\begin{itemize}
\item \textbf{Permutation invariance}: Attention is invariant to the order of items/candidates, making the network robust to different input sequences
\item \textbf{Variable-length sets}: Can process different numbers of items or candidates without architectural changes
\item \textbf{Learned relationships}: Automatically discovers which spatial and relational patterns lead to efficient packings
\end{itemize}

\subsubsection{Comparison with CNNs and MLPs}

\begin{table}[h]
\centering
\begin{tabular}{|l|p{5.5cm}|p{5.5cm}|}
\hline
\textbf{Architecture} & \textbf{Strengths} & \textbf{Weaknesses} \\
\hline
\textbf{MLP} & Simple, fast, good for fixed-size feature vectors & Cannot model spatial structure or variable sets \\
\hline
\textbf{CNN} & Excellent for grid-based representations (heightmaps) & Requires spatial discretization, limited relational reasoning \\
\hline
\textbf{Transformer} & Handles variable sets, learns relational patterns & Computationally expensive ($O(n^2)$ in sequence length), needs large datasets \\
\hline
\end{tabular}
\end{table}


\subsection{Set Transformers}

Standard Transformers are permutation-equivariant (output order matches input order) but produce one output per input element. For combinatorial optimization, we often need a \textbf{permutation-invariant} summary of a set, followed by set-to-element or set-to-set operations. \textbf{Set Transformers} \cite{LSK+19} extend Transformers with specialized attention mechanisms for set-structured data.

\subsubsection{Motivation}

In bin packing, we often need to:

\begin{enumerate}
\item \textbf{Aggregate item information}: Summarize all remaining items into a global context vector
\item \textbf{Score variable-size sets}: Evaluate sets of candidates of different sizes
\item \textbf{Maintain permutation invariance}: The set $\{\text{item}_1, \text{item}_2\}$ should be processed identically to $\{\text{item}_2, \text{item}_1\}$
\end{enumerate}

Set Transformers provide building blocks for these operations.

\subsubsection{Induced Set Attention Block (ISAB)}

Standard self-attention has $O(n^2)$ complexity for a set of size $n$. \textbf{ISAB} reduces this to $O(nm)$ by introducing $m$ \textbf{inducing points} $\mathbf{I} \in \mathbb{R}^{m \times d}$ (learnable parameters):

\begin{enumerate}
\item Aggregate input set $\mathbf{X}$ into inducing points:
   \begin{equation}
   \mathbf{H} = \text{Attention}(Q=\mathbf{I}, K=\mathbf{X}, V=\mathbf{X})
   \end{equation}

\item Project inducing points back to output:
   \begin{equation}
   \mathbf{Y} = \text{Attention}(Q=\mathbf{X}, K=\mathbf{H}, V=\mathbf{H})
   \end{equation}
\end{enumerate}

This is analogous to bottleneck layers in CNNs: compress information through inducing points, then expand.

\subsubsection{Pooling by Multihead Attention (PMA)}

To produce a \textbf{fixed-size} output from a variable-size set, PMA uses $k$ learnable \textbf{seed vectors} $\mathbf{S} \in \mathbb{R}^{k \times d}$:

\begin{equation}
\mathbf{Z} = \text{Attention}(Q=\mathbf{S}, K=\mathbf{X}, V=\mathbf{X})
\end{equation}

The seeds attend to the entire input set, producing $k$ summary vectors. For $k=1$, this gives a \textbf{single global representation} of the set, similar to global average pooling but learned.

\subsubsection{Application to Packing Problems}

\paragraph{Item Set Encoding:} Given a set of items $\{\mathbf{v}_i\}$ to be packed:

\begin{enumerate}
\item Embed each item: $\mathbf{e}_i = \text{Embed}(\mathbf{v}_i)$ capturing item properties
\item Apply ISAB layers: Learn inter-item relationships (e.g., grouping by size or shape compatibility)
\item Pool with PMA: Produce a global item set representation $\mathbf{z}_{\text{items}}$
\end{enumerate}

\paragraph{Candidate Scoring:} Given a set of candidate placements $\{\mathbf{p}_j\}$:

\begin{enumerate}
\item Embed each candidate: $\mathbf{c}_j = \text{Embed}(\mathbf{p}_j)$ capturing position and available space
\item Attend to item set:
   \begin{equation}
   \mathbf{h}_j = \text{Attention}(Q=\mathbf{c}_j, K=\mathbf{z}_{\text{items}}, V=\mathbf{z}_{\text{items}})
   \end{equation}
\item Score candidates: $Q(s, a_j) = \text{MLP}(\mathbf{h}_j)$
\end{enumerate}

\paragraph{Advantages:}

\begin{itemize}
\item \textbf{Handles variable-size inputs}: Number of items or candidates can vary
\item \textbf{Scalable}: ISAB reduces quadratic complexity to linear in the number of inducing points
\item \textbf{Interpretable attention}: Attention weights reveal which items influence which placements
\end{itemize}

\paragraph{Limitations:}

\begin{itemize}
\item \textbf{Inducing point tuning}: Choosing the number of inducing points $m$ requires experimentation
\item \textbf{Less common}: Fewer pre-trained models or established best practices compared to standard Transformers
\end{itemize}


\subsection{Deep Reinforcement Learning Context}

The neural architectures discussed above are typically integrated within \textbf{deep reinforcement learning (DRL)} frameworks, where the network serves as a function approximator for value functions or policies. This section briefly contextualizes the role of neural networks in the broader DRL pipeline.

\subsubsection{Value-Based vs. Policy-Based Methods}

\paragraph{Value-Based (DQN, Double DQN):}

\begin{itemize}
\item Learn $Q(s, a)$ or $V(s)$
\item Policy is derived implicitly: $\pi(s) = \arg\max_a Q(s, a)$
\item Off-policy: Can learn from experiences collected by different policies
\item Sample efficient but limited to discrete action spaces
\end{itemize}

\paragraph{Policy-Based (REINFORCE, PPO, A3C):}

\begin{itemize}
\item Learn $\pi(a|s; \theta)$ directly
\item Sample actions from the policy: $a \sim \pi(\cdot|s)$
\item On-policy: Requires recent experience from the current policy
\item Works with continuous action spaces but higher variance
\end{itemize}

\paragraph{Actor-Critic (A2C, PPO):}

\begin{itemize}
\item Combine both: Learn $\pi(a|s)$ (actor) and $V(s)$ or $Q(s, a)$ (critic)
\item Critic reduces variance of policy gradient estimates
\item State-of-the-art for many RL tasks
\end{itemize}

\subsubsection{Sample Efficiency and Curriculum Learning}

DRL often requires millions of environment steps to converge. For packing problems:

\begin{itemize}
\item \textbf{Simulated environments}: Fast simulation enables rapid data collection and policy evaluation
\item \textbf{Curriculum learning}: Train on easy instances (few items, regular shapes) before progressively harder ones
\item \textbf{Transfer learning}: Pre-train on synthetic problem distributions, fine-tune on target distributions
\end{itemize}


\subsection{Hybrid Architectures and Integration}

Modern packing systems often \textbf{combine} multiple neural components to leverage complementary strengths:

\paragraph{CNN + Transformer:}
\begin{itemize}
\item CNN processes spatial representations (heightmaps, occupancy grids) $\to$ spatial features
\item Transformer processes item and candidate embeddings $\to$ relational features
\item Concatenate and feed to decision head
\end{itemize}

\paragraph{GNN + DQN:}
\begin{itemize}
\item GNN encodes packing state as graph
\item DQN uses GNN embeddings to score actions
\item Experience replay and target networks stabilize training
\end{itemize}

\paragraph{Attention + Heuristics:}
\begin{itemize}
\item Neural network scores candidates
\item Heuristic filters infeasible placements
\item Combine scores: $\text{final\_score} = \alpha \cdot \text{NN\_score} + (1-\alpha) \cdot \text{heuristic\_score}$
\end{itemize}

This hybrid approach leverages the strengths of both learned and hand-crafted methods: neural networks adapt to data distributions, while heuristics guarantee feasibility and provide inductive biases.


\subsection{Summary}

Neural network architectures for 3D bin packing span multiple paradigms, each offering distinct capabilities for processing geometric, relational, and sequential aspects of the problem. This section reviewed foundational architectures before examining their integration with reinforcement learning frameworks.

\paragraph{Foundational Architectures:}

\begin{itemize}
\item \textbf{Convolutional Neural Networks (CNNs)}: Excel at processing spatial representations (heightmaps, occupancy grids, voxelized containers) through hierarchical feature learning. CNNs provide translation equivariance and parameter efficiency but require discretization and struggle with long-range relational reasoning.

\item \textbf{Graph Neural Networks (GNNs)}: Explicitly encode relational structure through message passing over graphs representing items, placements, and spatial relationships. GNNs naturally capture support dependencies and constraints but incur graph construction overhead and scalability challenges.

\item \textbf{Recurrent Neural Networks (RNNs)}: Model sequential dependencies in packing trajectories through internal hidden states. LSTM and GRU variants handle long-range temporal dependencies, making them suitable for tracking placement history and container evolution, though they process sequentially and cannot parallelize.
\end{itemize}

\paragraph{Reinforcement Learning Integration:}

These foundational architectures are typically integrated within deep reinforcement learning frameworks:

\begin{itemize}
\item \textbf{Deep Q-Networks (DQN)}: Learn value functions $Q(s, a)$ for placement decisions, using experience replay and target networks for stable training. DQN can incorporate any of the above architectures as its function approximator.

\item \textbf{Double DQN}: Addresses overestimation bias through decoupled action selection and evaluation, providing more stable learning in large action spaces.

\item \textbf{Transformers and Set Transformers}: Enable permutation-invariant relational reasoning over variable-size sets of items and candidates, learning attention patterns that identify compatible placements without manual feature engineering.
\end{itemize}

\paragraph{Architecture Selection Guidelines:}

The choice of architecture depends on problem characteristics and computational constraints:

\begin{itemize}
\item \textbf{Spatial grid representations}: CNN-based architectures for heightmaps and voxel grids
\item \textbf{Explicit relational constraints}: Graph Neural Networks for support and adjacency relationships
\item \textbf{Sequential placement trajectories}: RNN/LSTM for temporal dependency modeling
\item \textbf{Variable-size item/candidate sets}: Transformer or Set Transformer architectures
\item \textbf{Fixed-size feature vectors}: MLP-based networks within DQN framework
\end{itemize}

\paragraph{State-of-the-Art Approaches:}

Modern packing systems increasingly employ \textbf{hybrid architectures} that combine multiple paradigms to leverage complementary strengths:
\begin{itemize}
\item CNN for spatial feature extraction + Transformer for relational reasoning
\item GNN for constraint encoding + DQN for reinforcement learning
\item RNN for sequence modeling + Attention for candidate selection
\end{itemize}

These hybrid approaches balance expressiveness, computational cost, and sample efficiency, representing the current frontier in learned packing methods. The specific instantiation of these architectures for particular packing scenarios, including training procedures, hyperparameter choices, and domain-specific adaptations, is explored in subsequent implementation chapters.
